\noindent \textit{\textcolor{red!60!black}{
The authors made no comment about the fact that the 1933 and 1934 coefficients do not seem statistically different than the 1930 effects (Figure 3).
}}

\bigskip

\noindent The referee is right, and we cannot reject the null that the 1930 coefficient equals the 1933 or 1934 coefficient: with two-way clustered standard errors in a sample of this size, pairwise comparisons between individual years have little power. Two points nonetheless bear on the interpretation. First, the identifying contrast is between each litigation year and the omitted base year of 1932, not between the litigation years and 1930; relative to 1932, the 1933 and 1934 interactions are individually significant while no pre-treatment coefficient is. Second, the relevant question for parallel trends is whether the pre-treatment coefficients show significant departures from zero or a trend into the treatment window. Neither is present: the 1926--1931 coefficients in column 2 of Table 4 ($+0.040$, $+0.026$, $+0.012$, $+0.022$, $-0.032$, $-0.007$) are uniformly small and insignificant, and the two years immediately preceding the litigation sit essentially at zero.

\medskip
\noindent The 1930 dip ($-0.032$, insignificant) plausibly reflects the deepening Depression: gold-exposed firms were more leveraged (Table 2) and may have been somewhat more sensitive to the credit contraction. If that leverage sensitivity drove the 1933--1934 coefficients, however, it should reappear in the 1937--1938 recession, where we find no differential effect.

\medskip
\noindent Section 5 now acknowledges directly that the litigation-year coefficients cannot be statistically distinguished from individual pre-treatment estimates such as the 1930 coefficient, and that the support for the parallel-trends assumption rests on the full pre-treatment pattern.
